\documentclass[11pt,a4paper]{report}

\usepackage[T1]{fontenc}
\usepackage{lmodern}
\usepackage[a4paper,margin=2.6cm]{geometry}
\usepackage{amsmath,amssymb,amsthm}
\usepackage{bm}
\usepackage{booktabs}
\usepackage{graphicx}
\usepackage{xcolor}
\usepackage{tikz}
\usetikzlibrary{arrows.meta,positioning,calc,decorations.pathmorphing,fit,backgrounds}
\usepackage[most]{tcolorbox}
\usepackage[font=small,labelfont=bf]{caption}
\usepackage{enumitem}
\usepackage{microtype}
\usepackage{setspace}
\usepackage[numbers,sort&compress]{natbib}
\usepackage[colorlinks=true,linkcolor=black,citecolor=NavyBlue,urlcolor=NavyBlue]{hyperref}

\definecolor{NavyBlue}{RGB}{20,60,120}
\definecolor{DraftBlue}{RGB}{28,72,138}
\definecolor{DraftGrey}{RGB}{120,120,120}
\definecolor{PanelGrey}{RGB}{246,246,248}

\onehalfspacing
\setlength{\parskip}{0.25em}

% ---------------------------------------------------------------- theorems
\theoremstyle{plain}
\newtheorem{theorem}{Theorem}[chapter]
\newtheorem{proposition}[theorem]{Proposition}
\newtheorem{lemma}[theorem]{Lemma}
\theoremstyle{definition}
\newtheorem{definition}[theorem]{Definition}
\theoremstyle{remark}
\newtheorem{remark}[theorem]{Remark}

% ============================================================================
%  DRAFT ANNOTATION SYSTEM
%  Set \draftmodefalse to remove every annotation from the build.
% ============================================================================
\newif\ifdraftmode
\draftmodetrue

\ifdraftmode
  \newenvironment{reserved}[1]{%
    \begin{tcolorbox}[
      enhanced, breakable,
      colback=PanelGrey, colframe=DraftBlue,
      boxrule=0.5pt, left=6pt,right=6pt,top=5pt,bottom=5pt,
      fonttitle=\bfseries\small\sffamily,
      title={RESERVED \textendash\ #1}]
      \small\sffamily
  }{%
    \end{tcolorbox}
  }
\else
  \newenvironment{reserved}[1]{\setbox0=\vbox\bgroup}{\egroup}
\fi

\newcommand{\dnote}[2]{\ifdraftmode{\sffamily\small\color{DraftBlue}\textbf{[#1]}\ #2}\fi}

% #1 label  #2 caption  #3 specification of what the final figure must contain
\newcommand{\figplaceholder}[3]{%
  \begin{figure}[htbp]
  \centering
  \ifdraftmode
    \begin{tcolorbox}[
      width=0.92\textwidth, colback=PanelGrey, colframe=DraftGrey,
      boxrule=0.5pt, sharp corners, halign=flush left]
      \small\sffamily
      \textbf{\color{DraftBlue}FIGURE PLACEHOLDER}\\[3pt]
      #3
    \end{tcolorbox}
  \else
    \fbox{\parbox{0.9\textwidth}{\centering\vspace{2cm}[figure]\vspace{2cm}}}
  \fi
  \caption{#2}
  \label{#1}
  \end{figure}}

\newcommand{\tbd}{\ifdraftmode{\color{DraftBlue}\sffamily\textsc{tbd}}\else--\fi}

% ============================================================================
%  NOTATION
% ============================================================================
\newcommand{\R}{\mathbb{R}}
\newcommand{\Sph}{\mathbb{S}}
\newcommand{\SOthree}{\mathrm{SO}(3)}
\newcommand{\SEthree}{\mathrm{SE}(3)}
\newcommand{\sothree}{\mathfrak{so}(3)}
\newcommand{\Mfd}{\mathcal{M}}
\newcommand{\Exp}{\operatorname{Exp}}
\newcommand{\Log}{\operatorname{Log}}
\newcommand{\Ad}{\operatorname{Ad}}
\newcommand{\Tr}{\operatorname{tr}}
\newcommand{\E}{\mathbb{E}}
\newcommand{\Prob}{\mathbb{P}}
\newcommand{\dif}{\mathrm{d}}
\newcommand{\Loss}{\mathcal{L}}
\newcommand{\vfield}{u}
\newcommand{\pred}{v_\theta}
\newcommand{\Haar}{\mathcal{U}_{\SOthree}}
\newcommand{\ctx}{c}
\newcommand{\skewm}[1]{\widehat{#1}}
\newcommand{\Rot}{R}
\newcommand{\Ang}[1]{\ensuremath{#1^{\circ}}}
\newcommand{\divg}{\operatorname{div}}
\DeclareMathOperator*{\argmin}{arg\,min}

\title{\textbf{SigmaFlow}\\[0.3em]
       \large Riemannian Flow Matching for Generative Protein--Ligand Docking}
\author{Julian M\"uller}
\date{Draft of \today}

\begin{document}
\maketitle

% ---------------------------------------------------------------------------
\begin{abstract}
\noindent
Molecular docking predicts the bound geometry of a small molecule, the ligand,
in a protein binding site. A recent advance in this direction is SigmaDock
\citep{prat2026sigmadock}, which formulates docking as score-based Riemannian
diffusion over rigid molecular fragments in $\SEthree$. SigmaDock reaches a
Top-1 success rate of $79.9\%$ for poses that are simultaneously within
$2$\,\AA{} RMSD of the crystallographic pose and PoseBusters-valid, and is the
first deep-learning docking method to surpass classical physics-based docking
on the intended PoseBusters train-test split. This thesis replaces SigmaDock's
diffusion-based generative mechanism with Riemannian flow matching while
holding the network architecture, the data pipeline and the training budget
fixed, so that the effect of the generative mechanism is isolated from
confounding architectural or computational change. The resulting system,
SigmaFlow, provides, to our knowledge, the first flow-matching formulation of
protein-ligand docking under SigmaDock's fragment-based conditional
representation on $\SEthree^{N}$, and is evaluated against its diffusion
parent at matched training compute. Beyond that controlled comparison, the
thesis develops two variants that the flow formulation makes natural. The
first separates the translational and rotational readout into two heads on a
shared equivariant trunk, motivated by the observation that the inherited
readout asks a single per-atom field to serve two spatial moments at once. The
second attaches a confidence model trained on the generator's own samples, in
the manner of DiffDock \citep{corso2023diffdock}, which addresses the
observation that pose ranking rather than pose generation accounts for a large
part of the achievable accuracy. Throughout, the deterministic
ordinary-differential-equation formulation is exploited to study how accuracy
depends on the number of integration steps and on the number of independent
samples drawn.
\\\\
\dnote{RESULT PENDING}{The abstract closes with the verdict of the
matched-compute experiments and of the two variants. Rewrite the final two
sentences once Chapters~\ref{ch:minimal}, \ref{ch:separate} and
\ref{ch:confidence} carry their results.}
\end{abstract}

\tableofcontents

% ===========================================================================
\chapter{Introduction}
\label{ch:intro}

\section{Docking as a generative problem}

Molecular docking asks where a small molecule, the \emph{ligand}, sits when
bound to a protein, and in what orientation and internal conformation. The
answer determines whether a candidate compound is worth synthesising, and
computational prediction has been a central tool of structure-based drug
design for decades. Classical docking programs pose the question as
optimisation. They define a scoring function that approximates binding free
energy and then search the space of poses for its minimum. This approach has
important limitations. The scoring functions are only approximate surrogates
for the underlying molecular interactions, exploring the high-dimensional
conformational search space becomes computationally demanding in large-scale
virtual screening, and the resulting ranked candidate poses do not define an
explicit probability distribution over plausible binding configurations, which
makes uncertainty in the predicted pose difficult to quantify.

Generative modelling offers a different perspective. Rather than searching
directly for the optimum of a hand-designed scoring function, a generative
model learns a distribution over plausible binding configurations conditioned
on the protein, the ligand and the binding pocket. Docking then becomes the
task of sampling from
\begin{equation}
    p(\text{pose}\mid\text{protein},\text{ligand},\text{binding pocket}),
    \label{eq:docking-as-sampling}
\end{equation}
with multiple samples representing alternative hypotheses for how the ligand
may bind. DiffDock \citep{corso2023diffdock} demonstrated that diffusion-based
generative modelling can perform this task at scale, which motivated
subsequent work on more expressive molecular representations, on equivariant
architectures, and on diffusion processes adapted to the structure of the
docking problem.

A central development in that line is to respect the geometry of the pose
space itself. A molecular pose is not naturally represented by an
unconstrained Euclidean vector, since moving a rigid molecular fragment
requires a translation together with a rotation, and the pair forms an element
of the special Euclidean group $\SEthree$. Decomposing a flexible ligand into
$N$ rigid fragments therefore leads naturally to a configuration space of the
form $\SEthree^{N}$. SigmaDock \citep{prat2026sigmadock} exploits precisely
this representation, combining an equivariant architecture with a diffusion
process defined over fragment poses, and it reports that the fragment
representation is what makes the learning problem tractable. In parallel,
recent developments in generative modelling have introduced flow matching as
an alternative and broader framework for constructing continuous generative
processes \citep{lipman2023flow,liu2023rectified,albergo2023stochastic}. Its
close relationship to diffusion, together with the freedom to define direct
transport paths between a source and a target distribution, raises the
question of whether these developments can be exploited in the geometric
setting of molecular docking.

This thesis investigates that question through \emph{SigmaFlow}, which
replaces SigmaDock's diffusion dynamics with Riemannian flow matching while
retaining its fragment-based geometric representation and, as far as possible,
the surrounding model architecture. Instead of learning to reverse a
stochastic noising process, SigmaFlow learns a time-dependent vector field
that transports a simple reference distribution towards the conditional
distribution of bound poses through an ordinary differential equation on
$\SEthree^{N}$. This provides a controlled setting in which diffusion and flow
matching can be compared for molecular docking. The resulting continuous
normalising flow also admits evaluation of the density assigned to an
individual pose through the probability-flow dynamics, which provides a
model-based quantity that can be investigated as a measure of pose confidence
alongside the variability observed across generated samples.

\section{Research question and scope}

SigmaDock \citep{prat2026sigmadock} combines an $\SEthree$-equivariant
EquiformerV2 backbone \citep{liao2024equiformerv2} with a score-based
diffusion process over rigid-fragment poses, and it is the starting point of
this thesis. Diffusion is not the only way to learn a transport between a
tractable source distribution and the data. Flow matching
\citep{lipman2023flow} learns a velocity field directly by regression, without
simulating a stochastic process at training time, and its Riemannian extension
\citep{chen2024riemannian} carries the construction to manifolds. The two
approaches are closely related, since a diffusion model admits a deterministic
probability-flow ordinary differential equation \citep{song2021sde}, but they
differ in what the network is asked to predict, in how the source distribution
is chosen, and in how many network evaluations sampling requires. This thesis
asks:
\begin{quote}
\itshape
What changes when SigmaDock's score-based diffusion is replaced by Riemannian
flow matching on the protein-ligand pose manifold, while holding architecture,
data and computational budget fixed? Which flow-specific design choices are
needed to make this substitution competitive, and what new modelling
possibilities does it enable?
\end{quote}

\noindent
The design principle in the first part is \emph{minimal intervention}. The
backbone, the featurisation, the data pipeline, the optimiser and the
evaluation protocol are held methodologically fixed, adapted only where
technical integration requires it, and only the generative mechanism is
replaced. That is what makes the comparison interpretable, and
Chapter~\ref{ch:minimal} documents the extent to which it was achieved,
including where it was not. Two further variants then relax the constraint in
one controlled direction each. Chapter~\ref{ch:separate} separates the
translational and rotational readout into two heads on a shared trunk, and
Chapter~\ref{ch:confidence} adds a confidence model trained on the
generator's own samples.

Two limitations should be kept in view throughout. All experiments in this
thesis run at a fraction of the compute used to produce SigmaDock's published
numbers, so absolute performance is well below the literature and no benchmark
claim is made. What is claimed is a matched-budget comparison between
generative mechanisms. In addition, the evaluation set used here comprises $209$
complexes, whereas the published PoseBusters figure is computed over $308$.
The provenance of that difference is discussed in
Section~\ref{sec:evaluation-set}, and its consequence is that the two sets of
numbers are not directly comparable even at equal compute.

% ===========================================================================
\chapter{Docking as Conditional Generation on a Product Manifold}
\label{ch:problem}

This chapter states the problem that both generative mechanisms solve, in the
representation that SigmaDock introduced and that SigmaFlow inherits
unchanged. Section~\ref{sec:state-space} fixes the state space,
Section~\ref{sec:why-fragments} explains why the fragment representation is
the right one and cites the result that justifies it,
Section~\ref{sec:conditional} states the learning problem, and
Section~\ref{sec:sigmadock} describes the score-based baseline. The longest
part is Section~\ref{sec:backbone}, which traces exactly how the network
output becomes the two score components that enter the loss, because
Chapters~\ref{ch:minimal} and \ref{ch:separate} both turn on that chain.

\section{Poses, rigid fragments, and the state space}
\label{sec:state-space}

Let a ligand consist of $N_a$ atoms. A \emph{pose} is an assignment of
positions to those atoms. Modelling the full conformational space couples bond
lengths, bond angles and torsions. SigmaDock instead partitions the ligand
into $N$ rigid \emph{fragments}, treats internal fragment geometry as fixed at
a precomputed conformer, and generates only the rigid-body placement of each
fragment. Torsional flexibility survives implicitly, as the relative placement
of neighbouring fragments, but is not itself a modelled coordinate.

A single fragment's placement is an element of the special Euclidean group
\begin{equation}
  \SEthree \;=\; \bigl\{\, (\Rot, p) : \Rot \in \SOthree,\; p \in \R^3 \,\bigr\},
\end{equation}
acting on a body-frame atom offset $\tilde{x}$ by
$\tilde{x} \mapsto \Rot\,\tilde{x} + p$. The full state is therefore a point
on the product manifold
\begin{equation}
  \Mfd \;=\; \SEthree^{N}
  \;\cong\; \SOthree^{N} \times \bigl(\R^{3}\bigr)^{N},
  \qquad \dim \Mfd = 6N .
  \label{eq:state-space}
\end{equation}
Two conventions matter for what follows. The implementation treats $\Mfd$ as a
\emph{direct} product, so rotations and translations receive independent
probability paths and independent vector fields with no cross terms, which is
a modelling choice rather than a property of $\SEthree$. Translations are
expressed relative to the binding-pocket centroid, so the origin of the
$\R^{3}$ factor is a chemically meaningful point, and the model is not
required to be translation-equivariant.

The number of fragments is therefore the quantity that fixes the dimension of
the problem, and it varies considerably across the benchmark.
Figure~\ref{fig:fragment-histogram} reports its distribution. Over $208$ of the
$209$ evaluation complexes, one having exceeded the enumeration time limit,
the mean is $4.51$ fragments with a standard deviation of $2.11$, the median
is $4$, the mode is $5$, the interquartile range is $3$ to $6$, and the range
is $1$ to $11$. The distribution is right-skewed, $38.5\%$ of ligands have at
most three fragments, $18.8\%$ have at least seven, and five ligands have
exactly one, meaning that they carry no rotatable bond and are generated as a
single rigid body. The state dimension $D = 6N$ therefore averages $27.1$ and
never exceeds $66$. The space is low-dimensional, and the difficulty lies in
its geometry and in the conditioning rather than in its dimension.

\figplaceholder{fig:fragment-histogram}
{Distribution of the number of rigid fragments per ligand over the
PoseBusters evaluation set, with the induced state dimension $D = 6N$ on the
second axis.}
{Histogram of fragment count $N$ over the $208$ measured ligands, bars in a
single muted colour, with a cumulative curve overlaid on a secondary axis and
a second top axis labelled $D = 6N$. Annotate mean $4.51$, median $4$ and the
five single-fragment ligands. Source data:
\texttt{Thesis Visualisierungen/data/fragment\_distribution.csv}. This is the
figure that establishes the scale of the generative problem and should be
referred back to whenever fragment count is used as a covariate.}

\section{Why fragments, and what the representation buys}
\label{sec:why-fragments}

The alternative to fragments is the torsional parameterisation used by
DiffDock and by torsional diffusion more generally
\citep{corso2023diffdock,jing2022torsional}, in which the generative process
acts on the global rigid motion together with the $k$ dihedral angles. That
parameterisation is more compact, since it has $k+6$ degrees of freedom
against $6N$, but SigmaDock argues that it is worse conditioned, and states
the reason as a formal result.

\begin{theorem}[\citealp{prat2026sigmadock}, Theorem 1]
\label{thm:product-measure}
For standard molecular topologies, torsional models define nonlinear mappings
from torsion angles to Cartesian space, producing highly entangled,
non-product induced measures. In contrast, disjoint rigid fragments yield a
factorised product of Haar measures on $\SEthree^{m}$.
\end{theorem}

\noindent
The result is quoted, not reproved. Its consequence is what matters here. A
local change in a single torsion angle produces a non-local Cartesian
displacement, because every atom distal to the rotated bond moves, and the
displacement grows with the length of the chain through a lever effect. Under
the induced measure in Cartesian coordinates, where the network observes the
data, independent torsional perturbations therefore become correlated, and the
product structure is destroyed. Fragments avoid this because the forward
kernel factorises over disjoint $\SEthree$ components. Inter-fragment
correlation then enters only through the learnt score, not through the noise
model itself, which yields a better-conditioned learning problem and more
stable reverse-time integration.

A second difficulty with torsional parameterisations is a gauge ambiguity.
Mapping a torsional increment $\Delta\phi$ to a Cartesian displacement
requires a choice of which side of the rotatable bond moves, and
implementations resolve this with heuristics such as removing the induced net
rigid motion by superposition, which itself breaks the product structure.
SigmaDock avoids the question by construction.

Two further components of the representation are inherited unchanged and are
recorded here for completeness. First, the naive fragmentation at every
torsional bond yields $\hat{m} = k+1$ fragments, and SigmaDock reduces this by
a stochastic merging procedure that recursively merges adjacent fragments to
an irreducible set of size $m$, empirically about $\tfrac{2}{3}\hat{m}$, which
directly reduces the dimension of the generative problem. Second, geometric
information that fragmentation would otherwise discard is reintroduced as a
soft conditioning signal, by feeding relative cross-fragment distance
mismatches as edge features, which fixes bond lengths and bond angles across a
cut bond without constraining the dihedral itself
\citep[Lemma~1]{prat2026sigmadock}. Neither component is modified in this
thesis, and both apply identically to every arm of the comparison.

\section{The conditional generative formulation}
\label{sec:conditional}

Write $\ctx$ for the conditioning context, comprising the protein pocket, the
ligand's molecular graph and the precomputed fragment conformers. The learning
target is the conditional law of the crystallographic pose,
\begin{equation}
  p_{\mathrm{data}}(\,\cdot \mid \ctx\,) \quad\text{on}\quad \Mfd ,
\end{equation}
from which we wish to draw samples at inference. Two features distinguish this
from unconditional generative modelling on a manifold. The conditioning is
strong, in the sense that given a pocket the plausible poses occupy a small
region of $\Mfd$. And the training data provide exactly one sample per
context, namely the deposited structure, so there is no set of ground-truth
poses per complex to match in distribution. Every generative construction in
this thesis must therefore be conditional on $\ctx$ and trainable from single
observations.

\section{SigmaDock: the score-based baseline}
\label{sec:sigmadock}

SigmaDock defines a forward noising process that carries the data pose towards
a tractable reference. Writing $Z^{(t)} = (p^{(t)}, \Rot^{(t)})$ for the state
at diffusion time $t$, the forward process is the stochastic differential
equation
\begin{equation}
  \dif Z^{(t)}
  \;=\; \Bigl[-\tfrac12 p^{(t)},\, 0\Bigr]\dif t
   \;+\; \Bigl[\dif B^{(t)}_{\R^{m\times 3}},\, \dif B^{(t)}_{\SOthree^{m}}\Bigr],
  \label{eq:forward-sde}
\end{equation}
which is variance-preserving on the translational factor and Brownian motion
on the rotational factor \citep{yim2023se3,leach2022so3}. Its stationary
distribution is
\begin{equation}
  q(z) \;=\; \mathcal{N}(p;0,I) \otimes \Haar^{\otimes m},
  \label{eq:stationary}
\end{equation}
an isotropic Gaussian on the translations together with the uniform measure on
each rotational factor. This distribution will reappear in
Chapter~\ref{ch:theory} as SigmaFlow's source, which is a useful continuity
between the two mechanisms rather than a coincidence.

A network $s_\theta(z,t,\ctx)$ is trained by denoising score matching
\citep{vincent2011connection,song2019ncsn} to estimate the Stein score
$\nabla_z \log p_t$ of the perturbed marginals, understood as a Riemannian
gradient in the tangent space of $\SEthree^{m}$, and sampling integrates the
corresponding reverse-time stochastic differential equation
\citep{song2021sde}. Poses are generated in batches of $K$ independent seeds
and ranked by a fixed heuristic that combines a Vinardo pseudo binding energy
\citep{quiroga2016vinardo} with a set of physicochemical checks, and the
published headline numbers are post-ranking. SigmaDock deliberately does not
train a separate confidence model for this purpose, and
Chapter~\ref{ch:confidence} returns to that decision.

\subsection{The inherited backbone and how its output becomes a score}
\label{sec:backbone}

The network is an EquiformerV2 \citep{liao2024equiformerv2}, augmented by
SigmaDock with virtual nodes and edges over the chemical graph and with
structural-role-dependent featurisation. It is inherited by SigmaFlow without
modification, so this section states the interface rather than the internals,
and then traces the readout in full, because the readout is where the two
generative mechanisms diverge and where the variant of
Chapter~\ref{ch:separate} intervenes.

Physical correctness requires equivariance. Rotating the entire complex must
rotate the predicted pose update accordingly, and permuting atom indices must
induce the corresponding permutation of the outputs. A network built purely
from distances and other invariant scalars satisfies rotational invariance by
construction but cannot emit a directional update, because any function of
invariant quantities is itself invariant. The resolution is to propagate
\emph{steerable} features whose components transform according to prescribed
irreducible representations of $\SOthree$, indexed by degree $\ell$, with
$\ell=0$ components transforming as scalars and $\ell=1$ components as
vectors. Equivariant message passing then preserves these transformation laws
throughout the network. Permutation equivariance is automatic, because the
architecture is a graph neural network whose aggregation is permutation
invariant and whose outputs are attached to nodes.

The relevant consequence is concrete. The backbone maps the conditioning
context and the current noisy state to per-atom steerable features, and its
output block emits a single $\ell=1$ vector per node. Restricting to ligand
atoms and adopting the implementation's sign convention gives a per-atom field
\begin{equation}
  f_i \in \R^3, \qquad i = 1,\ldots,N_a .
  \label{eq:force-field}
\end{equation}
These vectors are not themselves scores. They form an intermediate atom-wise
equivariant representation which SigmaDock interprets as a \emph{pseudo-force
field}, in the sense that under a global rotation they transform in the same
way as physical force vectors, while being learned generative quantities
rather than forces derived from a molecular-mechanics potential.

The reason for reading the network output through mechanics rather than
predicting fragment updates directly is a gauge ambiguity in the fragment
representation, and it is worth stating because it explains an otherwise
surprising design. The global coordinates of a fragment are
$x_F = (p,\Rot)\cdot\tilde{x}_F$, but the body-frame coordinates $\tilde{x}_F$
have no canonical orientation. Replacing $\tilde{x}_F$ by
$\tilde{x}'_F = \Rot_0\tilde{x}_F$ and $(p,\Rot)$ by $(p, \Rot\Rot_0^{-1})$
describes the same physical fragment, so a head that predicted $\Rot$ directly
would have to commit to a convention. SigmaDock instead adapts the
$\SOthree$-equivariant prediction head of \citet{jin2023se3denoising}, based on
the Newton-Euler equations of rigid-body mechanics, and passes the backbone
output through it. With this choice the training objective and the sampling
procedure are invariant with respect to the orientation of local coordinates
\citep[Theorem~2]{prat2026sigmadock}.

Concretely, consider a fragment $F$ whose atoms occupy positions $r_i$ at
diffusion time $t$ and whose centre is $c_F$. The reduction produces a net
pseudo-force and a net pseudo-torque,
\begin{equation}
  F_F \;=\; \sum_{i \in F} f_i ,
  \qquad
  \tau_F \;=\; \sum_{i \in F} (r_i - c_F) \times f_i .
  \label{eq:newton-euler}
\end{equation}
The same atom-wise output is therefore read twice, in two different ways.
Translation depends on the zeroth spatial moment of the field over the
fragment, and rotation on its first moment about the fragment centre.

Since the coordinates are centred, $\sum_{i\in F}(r_i - c_F) = 0$, so a
spatially constant component of $f$ contributes to the net force but produces
no torque, while a spatially varying pattern whose atom-wise vectors sum to
zero contributes torque but no force. The two channels read genuinely
different spatial modes of one field. This observation is the entire
motivation for Chapter~\ref{ch:separate} and is worth carrying forward.

The reduction is completed with the fragment's mass and inertia. The
implementation assigns unit mass to every atom, so the fragment mass $M_F$ is
its atom count, and the inertia tensor about the centre is
\begin{equation}
  I_F \;=\; \sum_{i \in F}
  \Bigl[\, \|r_i - c_F\|_2^{2}\, I_3 \;-\; (r_i - c_F)(r_i - c_F)^{\top} \Bigr],
  \label{eq:inertia}
\end{equation}
regularised by clamping small eigenvalues before inversion. The Newton-Euler
step then converts force and torque into rigid-body increments,
\begin{equation}
  \Delta T_F \;=\; \frac{F_F}{M_F} \;=\; \frac{1}{|F|}\sum_{i\in F} f_i,
  \qquad
  \omega_F \;=\; I_F^{-1}\tau_F ,
  \qquad
  \skewm{\omega}_F \;=\; \operatorname{hat}(\omega_F).
  \label{eq:updates}
\end{equation}
Because the masses are uniform, the translational increment is exactly the
\emph{mean} of the per-atom vectors over the fragment. The rotational
increment carries an inverse inertia, which accounts for the fact that the
same torque need not produce the same rotational response for fragments with
different spatial distributions of atoms, and which is also, as
Chapter~\ref{ch:separate} discusses, the numerically delicate part of the
chain.

Only at this point does the diffusion parameterisation enter. Both increments
are rescaled by the score scaling of the corresponding diffuser at time $t$.
On the translational factor, with $\alpha_t = \exp(-\tfrac12\int_0^t \beta_s\,\dif s)$
and conditional standard deviation $\sigma_t = \sqrt{1-\alpha_t^{2}}$, the
scaling is $\lambda_{\R^3}(t) = 1/\sigma_t$ and the predicted score is
\begin{equation}
  s^{p}_\theta \;=\; \lambda_{\R^3}(t)\, \Delta T_F
  \;=\; \frac{1}{\sigma_t}\,\frac{1}{|F|}\sum_{i\in F} f_i .
  \label{eq:trans-score}
\end{equation}
This gives the readout the conditioning of an $\epsilon$-prediction
parameterisation, since the conditional score of the variance-preserving path
is $\nabla \log p_t(x_t\mid x_0) = -(x_t - \alpha_t x_0)/\sigma_t^{2}$ and the
network is thereby asked to predict a quantity of roughly unit scale at every
noise level.

On the rotational factor the production configuration takes a route that is
worth stating exactly, because it is easy to describe incorrectly. Rather than
emitting a score directly, the model first forms an estimate of the clean
rotation. Writing $\lambda_{\SOthree}(t)$ for the tabulated root-mean-square
scaling of the conditional $\SOthree$ score at time $t$, the increment is
exponentiated and applied to the current rotation,
\begin{equation}
  \Delta \Rot_F \;=\; \Exp\!\bigl(\lambda_{\SOthree}(t)\,\skewm{\omega}_F\bigr),
  \qquad
  \widehat{\Rot}_0 \;=\; \Delta \Rot_F\, \Rot_t ,
  \label{eq:rot-denoised}
\end{equation}
and the score is then evaluated in closed form from the pair
$(\Rot_t, \widehat{\Rot}_0)$ using the analytic conditional score of the
isotropic Gaussian on $\SOthree$,
\begin{equation}
  s^{\Rot}_\theta
  \;=\; \Rot_t \,\Log\bigl(\widehat{\Rot}_0^{\top}\Rot_t\bigr)\,
        \frac{\partial_\Omega \log f\bigl(\Omega,\sigma_t\bigr)}{\Omega},
  \qquad
  \Omega = \bigl\|\Log\bigl(\widehat{\Rot}_0^{\top}\Rot_t\bigr)\bigr\| ,
  \label{eq:rot-score}
\end{equation}
where $f(\,\cdot\,,\sigma)$ is the radial density of the $\mathrm{IGSO}(3)$
distribution. The network therefore predicts a \emph{denoised rotation} and
the score is a deterministic function of it. Two points follow. The rotational
head is a denoiser rather than a score head, which is precisely the structure
that makes the flow-matching substitution of Chapter~\ref{ch:theory} natural,
since flow matching also predicts something about the endpoint. And the score
in \eqref{eq:rot-score} is written as $\Rot_t X$ with $X \in \sothree$, which
fixes the trivialisation convention that Section~\ref{sec:trivialisation}
makes explicit.

The complete inherited mapping is therefore
\begin{equation}
  \{f_i\}_{i\in F}
  \;\longrightarrow\;
  \begin{cases}
    F_F \;\longrightarrow\; \Delta T_F \;\longrightarrow\; s^{p}_\theta, \\[0.6ex]
    \tau_F \;\longrightarrow\; I_F^{-1}\tau_F \;\longrightarrow\;
      \widehat{\Rot}_0 \;\longrightarrow\; s^{\Rot}_\theta .
  \end{cases}
  \label{eq:readout-summary}
\end{equation}
There is no independently learned translational head and rotational head at
the level of the backbone. Both score components originate in the same
atom-wise equivariant field, and the transformations between that field and
the two scores are substantially different.

The training objective regresses both predicted scores onto their analytic
conditional counterparts, weighted by the inverse squared score scaling of the
respective factor, and adds two reconstruction terms in data space, one on the
fragment centre and one geodesic term of the form
$\|\Log(\widehat{\Rot}_0^{\top}\Rot_0)\|_F^{2}$ on the fragment rotation. The
weights are inherited unchanged by every arm of the comparison in this thesis.

\figplaceholder{fig:architecture}
{The inherited architecture and its interface to the generative mechanism.
Components in grey are shared by every arm; only the shaded block is replaced.}
{Schematic in three columns.
\textbf{Left, shared:} protein pocket and ligand graph, featurisation,
EquiformerV2 backbone, per-atom steerable features.
\textbf{Middle, shared:} the $\ell=1$ output block emitting $f_i \in \R^3$
per atom, and the Newton-Euler reduction to $(F_F, \tau_F)$ per fragment,
Eq.~\eqref{eq:newton-euler}, with the mass and inertia step
Eq.~\eqref{eq:updates}. Draw the zeroth and first spatial moment explicitly,
so that the reader sees the two channels reading the same field differently.
\textbf{Right, replaced:} the interpretation layer. SigmaDock reads the
increments as a \emph{score}, rescales by $\lambda(t)$, and integrates a
reverse SDE; SigmaFlow reads them as a \emph{velocity} and integrates an ODE.
The figure must make visually obvious that the substitution is confined to one
block, since that is the evidence for the controlled comparison, and it
replaces roughly $400$ words of prose.}

\section{Evaluation}
\label{sec:metrics}

Docking accuracy is reported as the symmetry-corrected root-mean-square
deviation between predicted and crystallographic ligand atoms
\citep{meli2020spyrmsd}, computed against the nearest crystallographic copy,
with the fraction of complexes below $2$\,\AA{} as the conventional success
criterion. Physical plausibility is assessed with the PoseBusters battery
\citep{buttenschoen2024posebusters}. The primary criterion throughout is the
conjunction
\begin{equation}
  \mathrm{RMSD} < 2\,\text{\AA}
  \qquad\text{and}\qquad
  \text{PoseBusters-valid},
  \label{eq:success}
\end{equation}
because a geometrically close prediction achieved through a chemically invalid
configuration is not a successful prediction. Because RMSD mixes placement
error with internal geometry error, RMSD after optimal rigid superposition
\citep{kabsch1976solution} and the centroid displacement are reported
alongside it, since they isolate the two contributions.

Because sampling is repeated, a success rate is only defined relative to a
sampling and ranking protocol. Two quantities are distinguished throughout.
The \emph{oracle} rate at $K$ samples is the fraction of complexes for which at
least one of $K$ generated poses satisfies \eqref{eq:success}, and it is an
upper bound attainable only with a perfect ranker. The \emph{Top-1} rate is the
fraction for which the pose selected by the ranking heuristic satisfies it.
The gap between the two measures how much accuracy is lost in ranking rather
than in generation, and Chapter~\ref{ch:confidence} is devoted to it.

For the rotational channel a null model is needed, because the natural scale of
an orientation error is not obvious. If fragment orientations were drawn from
the Haar measure, the expected geodesic angle to the truth would be
\begin{equation}
  \E[\theta] \;=\; \frac{\pi^{2}/2 + 2}{\pi} \;\approx\; \Ang{126.5},
  \label{eq:haar-null}
\end{equation}
with a median of about $\Ang{132.3}$ under the sampling protocol used here.
Every rotational number reported in this thesis is stated against that
reference, because an error of $\Ang{130}$ is not a poor prediction but no
prediction at all.

\subsection{The evaluation set, and how it differs from the published one}
\label{sec:evaluation-set}

The published SigmaDock results are computed over the PoseBusters v2 set,
described as a temporal split of $308$ protein-ligand complexes whose protein
sequences are unseen and whose structures were released from 2021 onwards
\citep{prat2026sigmadock}, with a further $85$ complexes from the Astex diverse
set \citep{hartshorn2007astex} as a second evaluation set. The original
PoseBusters benchmark from which it derives contains $428$ complexes
\citep{buttenschoen2024posebusters}.

The copy of the benchmark available to this project contains $209$ complexes,
and this section records what is and is not known about that difference,
because it bounds what any comparison against the literature can claim.

Three checks establish that the reduction is a property of the data copy rather
than a silent filter in the pipeline. The benchmark directory contains $209$
subdirectories and yields $209$ recognised protein-ligand pairs, so the loader
discards nothing. The loader itself applies no selection beyond requiring a
matching structure file and ligand file per directory, and it emits a warning
whenever the two counts disagree, which it does not here. And the evaluation
stage, which globs the reference ligand files independently of the loader, also
finds $209$. The number is therefore fixed upstream of any modelling decision
taken in this thesis.

How the reduction from the published set came about is not documented and was
not reconstructible. The $209$ complexes carry $209$ distinct PDB identifiers
and $209$ distinct ligand codes, so the set contains no repeated protein and no
repeated ligand, and about ninety per cent of the identifiers fall in the
ranges assigned to recent depositions, which is consistent with but does not
establish a temporal split. Whether the $209$ form a subset of the published
$308$ has not been verified.

Two consequences follow, and they are different in kind. No absolute number in
this thesis may be compared against a published PoseBusters figure, since the
sets differ in size and in unknown composition, and this holds even for the
reproduced SigmaDock baseline. The internal comparison, by contrast, is
unaffected, because every arm is trained on the same data and evaluated on the
same $209$ complexes, and it is the internal comparison that this thesis
claims.

\dnote{OPEN}{Settling whether the $209$ are a subset of the published $308$
requires only the official identifier list. If they are a strict subset the
matter is closed with one sentence; if they are not, the composition of the
evaluation set needs to be characterised before any external comparison is
attempted.}

% ===========================================================================
\chapter{From Score to Velocity: Riemannian Flow Matching}
\label{ch:theory}

This chapter develops the generative machinery that replaces diffusion in
SigmaFlow. The exposition follows the development in
\citet{holderrieth2025flowmatching}, and results that are standard are stated
with attribution rather than reproved. What is derived here is only what the
later chapters depend on in a specific form, namely the construction on
$\SOthree$ and the trivialisation convention that the implementation must fix.

\section{Generation as sampling, and flows as the machinery}

Following \citet{holderrieth2025flowmatching}, generation is modelled as
sampling. The object to be generated is a point $z$ of the state space, the
data are a finite collection of such points drawn from an unknown distribution
$p_{\mathrm{data}}$, and a generative model is an algorithm that returns
further samples from it. In the conditional case, which is the relevant one
here, the target is $p_{\mathrm{data}}(\cdot\mid\ctx)$ as in
\eqref{eq:docking-as-sampling}.

A flow model produces such a sample by drawing an initial state from a
tractable source distribution and transporting it along a time-dependent
vector field. The trajectory solves the ordinary differential equation
\begin{equation}
  \frac{\dif}{\dif t} X_t \;=\; \vfield_t(X_t),
  \qquad X_0 \sim p_{\mathrm{init}},
  \label{eq:flow-ode}
\end{equation}
and the endpoint $X_1$ is the generated sample. The generative goal is
therefore to find a vector field whose flow satisfies
$X_0 \sim p_{\mathrm{init}} \Rightarrow X_1 \sim p_{\mathrm{data}}$. Two
questions follow, namely which vector field has that property and how it can
be learned from data.

\section{Conditional paths and the marginal vector field}

Flow matching answers the first question by specifying how probability mass
should move. For a data point $z \sim p_{\mathrm{data}}$ one chooses a
\emph{conditional probability path} $p_t(\cdot \mid z)$ with boundary
conditions
\begin{equation}
  p_0(\cdot \mid z) = p_{\mathrm{init}},
  \qquad
  p_1(\cdot \mid z) = \delta_z ,
  \label{eq:path-boundaries}
\end{equation}
so that conditional on a particular training example the path transports the
source continuously onto that example. Averaging over the data induces the
\emph{marginal probability path}
\begin{equation}
  p_t(x) \;=\; \int p_t(x\mid z)\, p_{\mathrm{data}}(z)\, \dif z ,
  \label{eq:marginal-path}
\end{equation}
which by construction interpolates from $p_{\mathrm{init}}$ at $t=0$ to
$p_{\mathrm{data}}$ at $t=1$.

If for each $z$ there is a tractable conditional vector field
$\vfield_t^{\mathrm{tgt}}(x\mid z)$ whose flow generates
$p_t(\cdot\mid z)$, then the field
\begin{equation}
  \vfield_t^{\mathrm{tgt}}(x)
  \;=\; \int \vfield_t^{\mathrm{tgt}}(x\mid z)\,
        \frac{p_t(x\mid z)\,p_{\mathrm{data}}(z)}{p_t(x)}\,\dif z
  \label{eq:marginal-field}
\end{equation}
generates the marginal path $p_t$, in the sense that it satisfies the
continuity equation for $p_t$
\citep[Theorem~9]{holderrieth2025flowmatching}. Integrating
\eqref{eq:flow-ode} under this field therefore transports
$p_{\mathrm{init}}$ to $p_{\mathrm{data}}$, which answers the first question.

\section{The conditional flow matching objective}

The field \eqref{eq:marginal-field} is intractable, since it averages against
the unknown data distribution. The second question is answered by not
evaluating it. One trains a neural field $\pred$ against the tractable
conditional field instead, through the conditional flow matching objective
\begin{equation}
  \Loss_{\mathrm{CFM}}(\theta)
  \;=\; \E_{\,t\sim\mathcal{U}[0,1],\; z\sim p_{\mathrm{data}},\;
             X_t \sim p_t(\cdot\mid z)}
    \Bigl[\bigl\| \pred(t, X_t) - \vfield_t^{\mathrm{tgt}}(X_t\mid z)\bigr\|^{2}\Bigr].
  \label{eq:cfm-loss}
\end{equation}
That this is legitimate is the central result of flow matching.

\begin{theorem}[\citealp{lipman2023flow}; \citealp{holderrieth2025flowmatching}, Theorem~12]
\label{thm:cfm}
The marginal flow matching loss and the conditional flow matching loss differ
by a constant independent of $\theta$, so
$\nabla_\theta \Loss_{\mathrm{FM}}(\theta)
 = \nabla_\theta \Loss_{\mathrm{CFM}}(\theta)$.
In particular, for the minimiser $\theta^{*}$ of $\Loss_{\mathrm{CFM}}$ and a
sufficiently expressive parameterisation, $\pred[\theta^{*}]$ equals the
marginal vector field $\vfield_t^{\mathrm{tgt}}$.
\end{theorem}

\noindent
By explicitly regressing against a tractable conditional field, one implicitly
regresses against the intractable marginal field that is actually wanted. Two
practical consequences follow and both matter for this thesis. Training is
\emph{simulation-free}, since no ODE is integrated during training and the
objective is an ordinary regression, and the probability path becomes an
explicit design choice rather than a consequence of a noising process.

In $\R^{d}$ the standard choice is the straight-line path between a source
draw and the data point,
\begin{equation}
  x_t = (1-t)\,x_0 + t\,x_1,
  \qquad
  \vfield_t(x_t \mid x_0, x_1) = x_1 - x_0 ,
  \label{eq:linear-path}
\end{equation}
whose conditional trajectories are straight and whose target is constant along
each trajectory. This is what SigmaFlow uses on the translational factor.

\section{Score and velocity as two coordinates on one object}
\label{sec:score-vs-velocity}

The title of this chapter is meant literally, and it is worth making the
relationship precise, since the substitution studied here is often described
loosely. For a Gaussian conditional path
$p_t(\cdot\mid z) = \mathcal{N}(\alpha_t z, \beta_t^{2} I_d)$ the marginal
velocity and the marginal score determine one another through an affine
relation,
\begin{equation}
  \vfield_t^{\mathrm{tgt}}(x)
  \;=\; a_t\, \nabla \log p_t(x) \;+\; b_t\, x,
  \qquad
  a_t = \beta_t^{2}\frac{\dot\alpha_t}{\alpha_t} - \dot\beta_t \beta_t,
  \qquad
  b_t = \frac{\dot\alpha_t}{\alpha_t},
  \label{eq:score-velocity}
\end{equation}
so on such a path the two objects carry the same information
\citep[Proposition~1]{holderrieth2025flowmatching}. Moreover, adding the score
to the deterministic dynamics with any diffusion coefficient $\sigma_t \ge 0$
preserves the marginals,
\begin{equation}
  \dif X_t \;=\; \Bigl[\vfield_t^{\mathrm{tgt}}(X_t)
    + \tfrac{\sigma_t^{2}}{2}\nabla \log p_t(X_t)\Bigr]\dif t
    \;+\; \sigma_t\, \dif W_t
  \;\Longrightarrow\; X_t \sim p_t ,
  \label{eq:sde-family}
\end{equation}
which recovers the deterministic probability-flow ODE at $\sigma_t = 0$ and a
diffusion sampler at $\sigma_t > 0$
\citep[Theorem~17]{holderrieth2025flowmatching}; the same statement in the
diffusion literature is the probability-flow ODE of \citet{song2021sde}.

Two things follow for the present work. First, diffusion and flow matching are
not competing theories but different coordinates on a family of processes
sharing the same marginals, and the substitution studied in this thesis is a
change of parameterisation together with a change of probability path.
Diffusion may be viewed as flow matching along a particular Gaussian path
\citep{lipman2023flow}. Second, and this is the practical content, the
relation \eqref{eq:score-velocity} is time-dependent through $a_t$ and $b_t$.
A score must be rescaled by noise-level-dependent factors before it becomes a
displacement, which is exactly what \eqref{eq:trans-score} and
\eqref{eq:rot-denoised} do, whereas a velocity is already in the units of the
state. The entire scalar apparatus that converts scores into displacements is
therefore deleted rather than modified when the substitution is made, and
Chapter~\ref{ch:minimal} records that this is what was done.

\section{The Riemannian generalisation}
\label{sec:riemannian}

On a Riemannian manifold $(\Mfd, g)$ several ingredients of the Euclidean
construction stop making sense, and it is worth being precise about which.
The interpolant $(1-t)x_0 + t x_1$ is undefined, because there is no addition.
The difference $x_1 - x_0$ is undefined. A velocity at $x$ is an element of the
tangent space $T_x\Mfd$, and tangent spaces at different points are different
vector spaces that cannot be compared without extra structure. Finally, the
squared error in \eqref{eq:cfm-loss} must be measured in the metric at the
point where both vectors live.

What survives is what matters. The argument behind Theorem~\ref{thm:cfm} uses
only bilinearity of an inner product and the tower property of conditional
expectation, both of which hold pointwise on a manifold, so the conditional
objective remains a valid surrogate \citep{chen2024riemannian}. Three
substitutions carry the construction over. Straight lines become geodesics,
differences become logarithmic maps, and the Euclidean norm becomes the
Riemannian metric at the relevant point,
\begin{equation}
  \Loss_{\mathrm{RCFM}}(\theta)
  \;=\; \E_{t,\,x_1,\,x_t}\bigl\|\, \pred(t,x_t)
        - \vfield_t(x_t \mid x_1)\,\bigr\|^{2}_{g(x_t)} .
  \label{eq:rcfm-loss}
\end{equation}
The geodesic interpolant between $x_0$ and $x_1$ is
\begin{equation}
  x_t \;=\; \Exp_{x_0}\!\bigl(t\, \Log_{x_0}(x_1)\bigr),
  \label{eq:geodesic-interp}
\end{equation}
where $\Exp_x : T_x\Mfd \to \Mfd$ follows the geodesic leaving $x$ with a given
initial velocity and $\Log_x$ is its local inverse. Differentiating
\eqref{eq:geodesic-interp} and re-expressing the result at $x_t$ gives the
conditional field
\begin{equation}
  \vfield_t(x_t \mid x_1) \;=\; \frac{\Log_{x_t}(x_1)}{1-t},
  \label{eq:riem-target}
\end{equation}
which states that at time $t$ one must cover the remaining geodesic distance
in the remaining time. On a simple geometry such as a compact Lie group with a
bi-invariant metric, where $\Exp$ and $\Log$ are available in closed form, this
construction remains simulation-free \citep{chen2024riemannian}.

\section{The rotation group and the trivialisation convention}
\label{sec:so3}

$\SOthree = \{\Rot\in\R^{3\times3} : \Rot^\top\Rot = I,\ \det\Rot = 1\}$ is a
compact three-dimensional Lie group. Its Lie algebra $\sothree$ is the space of
skew-symmetric matrices, identified with $\R^3$ by the hat map
\begin{equation}
  \skewm{\omega} =
  \begin{pmatrix}
    0 & -\omega_3 & \omega_2\\
    \omega_3 & 0 & -\omega_1\\
    -\omega_2 & \omega_1 & 0
  \end{pmatrix},
  \qquad \omega\in\R^3 .
\end{equation}
With the bi-invariant metric $\langle A,B\rangle = \tfrac12\Tr(A^\top B)$ the
exponential map is Rodrigues' formula and its inverse is available in closed
form away from the cut locus,
\begin{equation}
  \Exp(\skewm{\omega}) = I + \frac{\sin\theta}{\theta}\skewm{\omega}
    + \frac{1-\cos\theta}{\theta^{2}}\skewm{\omega}^{2},
  \qquad \theta = \|\omega\| ,
  \label{eq:rodrigues}
\end{equation}
\begin{equation}
  \Log(\Rot) = \frac{\theta}{2\sin\theta}\bigl(\Rot - \Rot^\top\bigr),
  \qquad
  \theta = \arccos\!\Bigl(\frac{\Tr \Rot - 1}{2}\Bigr).
  \label{eq:so3-log}
\end{equation}
The choice of a bi-invariant metric is not cosmetic. It makes left and right
translation isometries, so distances are unchanged by re-expressing both
arguments in another frame. It makes the geodesics through the identity
exactly the one-parameter subgroups $t \mapsto \Exp(t\skewm{\omega})$, which is
what puts \eqref{eq:rodrigues} in closed form. And it makes the Haar measure
the Riemannian volume, so the phrase ``uniform on $\SOthree$'' is unambiguous.
Every closed form in this chapter depends on it, and it is available because
$\SOthree$ is compact.

The geodesic distance between $\Rot_a$ and $\Rot_b$ is
$\|\Log(\Rot_a^\top\Rot_b)\|$, taking values in $[0,\pi]$. The logarithm is
branch-ambiguous at $\theta = \pi$, the cut locus, since a rotation by exactly
$\pi$ about an axis is equally a rotation by $-\pi$ about the same axis. This
is not a rare edge case in the present construction. Under the Haar measure
the angle has density proportional to $1 - \cos\theta$, which is increasing on
$[0,\pi]$, so uniformly sampled rotations concentrate towards large angles and
roughly one draw in ten exceeds $\Ang{170}$. Since SigmaFlow's rotational
source is exactly Haar, a tenth of all training pairs begin near the cut locus,
and the numerical treatment of \eqref{eq:so3-log} there is a correctness
requirement rather than a refinement.

\subsection{Body frame or spatial frame, fixed explicitly}
\label{sec:trivialisation}

Because $\SOthree$ is a group, a tangent vector at $\Rot$ may be transported to
the identity in two inequivalent ways, and the resulting conventions are easy
to conflate. Writing a curve's velocity as
\begin{equation}
  \dot{\Rot} = \Rot\,\skewm{\Omega}
  \;\Longleftrightarrow\;
  \skewm{\Omega} = \Rot^\top \dot{\Rot}
  \qquad \text{(body frame)}
  \label{eq:body-frame}
\end{equation}
differs from
\begin{equation}
  \dot{\Rot} = \skewm{w}\,\Rot
  \;\Longleftrightarrow\;
  \skewm{w} = \dot{\Rot}\,\Rot^\top
  \qquad \text{(spatial frame)},
  \label{eq:spatial-frame}
\end{equation}
and the two are related by the adjoint action,
\begin{equation}
  \skewm{\Omega} \;=\; \Rot^\top \skewm{w}\,\Rot
  \;=\; \Ad_{\Rot^{-1}}\skewm{w} .
  \label{eq:adjoint}
\end{equation}
This thesis uses the body-frame convention \eqref{eq:body-frame} throughout,
following \citet{marsden1999introduction}. Every rotational velocity, every
target and every network output is a body-frame quantity, and the integrator
correspondingly right-multiplies by the exponential.

Fixing this explicitly is necessary rather than pedantic, for a reason that is
visible in the baseline itself. The inherited rotational increment
$\skewm{\omega}_F$ of \eqref{eq:updates} is built from torques computed in the
global frame, so it is a spatial-frame quantity, and SigmaDock consumes it
consistently by left-multiplying in \eqref{eq:rot-denoised}. A body-frame
formulation cannot use it unconverted. Section~\ref{sec:sf-readout} states the
conversion that SigmaFlow applies.

\section{The SigmaFlow probability path and objective}
\label{sec:sigmaflow-objective}

Specialising \eqref{eq:geodesic-interp} and \eqref{eq:riem-target} to the two
factors of $\Mfd$ gives the construction implemented in SigmaFlow. The source
is the same distribution that SigmaDock's forward process converges to in
\eqref{eq:stationary}, namely an isotropic Gaussian on each translational
factor in pocket-centred coordinates and the Haar measure on each rotational
factor, so the two mechanisms start their sampling from the same law.

On each translational factor, with $x_0 \sim \mathcal{N}(0,I_3)$,
\begin{equation}
  x_t = (1-t)x_0 + t\,x_1,
  \qquad
  \vfield_t^{\mathrm{trans}} = x_1 - x_0 \;=\; \frac{x_1 - x_t}{1-t}.
  \label{eq:sf-trans}
\end{equation}
On each rotational factor, with $\Rot_0 \sim \Haar$,
\begin{equation}
  \Rot_t \;=\; \Rot_0\, \Exp\!\bigl(t\, \Log(\Rot_0^\top \Rot_1)\bigr),
  \qquad
  \vfield_t^{\mathrm{rot}} \;=\; \Log(\Rot_0^\top\Rot_1)
  \;=\; \frac{\Log(\Rot_t^\top \Rot_1)}{1-t}.
  \label{eq:sf-rot}
\end{equation}
In both cases the two expressions for the target agree. The right-hand forms,
which involve the current state, express the target as the remaining distance
divided by the remaining time and are singular at $t=1$; the left-hand forms,
which involve only the endpoints, are constant along each conditional
trajectory and are not singular. The implementation uses the endpoint forms
for exactly that reason, and the equivalent state-dependent forms are retained
for diagnostics.

The product structure of $\Mfd$ makes the extension to $N$ fragments
immediate. On a Riemannian product the metric is the direct sum of the factor
metrics, geodesics are products of factor geodesics, and $\Exp$ and $\Log$ act
componentwise, so a path constructed independently on each factor is a
legitimate path on $\Mfd$ and the conditional field is the concatenation of the
factor fields. The objective is therefore a weighted sum of the two factor-wise
losses,
\begin{equation}
  \Loss(\theta) \;=\;
  \E_{t\sim\mathcal{U}[0,1]}\Bigl[
    \lambda_{\mathrm{trans}}\bigl\|\pred^{\mathrm{trans}}
      - \vfield_t^{\mathrm{trans}}\bigr\|^{2}
    + \lambda_{\mathrm{rot}}\bigl\|\pred^{\mathrm{rot}}
      - \vfield_t^{\mathrm{rot}}\bigr\|_F^{2}\Bigr],
  \label{eq:sf-objective}
\end{equation}
with the rotational term measured in the Frobenius norm on $\sothree$ and the
weights inherited unchanged from SigmaDock's production configuration.
Sampling integrates the learned field by explicit Euler, using the exponential
map as the retraction on the rotational factor,
\begin{equation}
  x^{(k+1)} = x^{(k)} + \pred^{\mathrm{trans}}\Delta t,
  \qquad
  \Rot^{(k+1)} = \Rot^{(k)}\Exp\!\bigl(\pred^{\mathrm{rot}}\Delta t\bigr).
  \label{eq:euler}
\end{equation}
The right-multiplication in \eqref{eq:euler} is the concrete expression of the
body-frame convention fixed in Section~\ref{sec:trivialisation}.

\figplaceholder{fig:paths}
{Conditional probability paths on the two factors of $\Mfd$, with the
conditional target drawn at three times.}
{Two panels.
\textbf{Left:} $\R^3$ drawn as $\R^2$. A source sample $x_0$ from a Gaussian
centred at the pocket origin, the data point $x_1$, the straight interpolant,
and the constant target $x_1-x_0$ drawn as equal arrows at
$t = 0.2, 0.5, 0.8$.
\textbf{Right:} $\SOthree$ drawn as a schematic curved surface. $\Rot_0$
sampled Haar-uniformly, $\Rot_1$ the crystal orientation, the geodesic between
them, and the body-frame target drawn at the same three times, with identical
length at every $t$. The visual point is that neither set of arrows shrinks as
$t \to 1$, which is what distinguishes the endpoint parameterisation from the
state-dependent one. Replaces roughly $250$ words.}

\section{The source distribution as a modelling choice}
\label{sec:source-theory}

One consequence of the freedom granted by Theorem~\ref{thm:cfm} deserves
emphasis, because Section~\ref{sec:future-source} returns to it. Under
diffusion the source is not chosen. It is whatever the forward noising process
converges to, which is \eqref{eq:stationary}. Under flow matching the only
requirements on $p_0$ are that it be tractable to sample and that the
conditional path connect it to $p_1$. Any such $p_0$ is admissible, including
one that depends on the conditioning context $\ctx$.

This is a genuinely flow-specific degree of freedom, and it is not merely a
convenience, because the difficulty of the learning problem depends on how much
the source already knows. SigmaFlow's translational source is centred on the
binding pocket and is therefore informative, in the sense that its mean is the
right order of magnitude from the answer. Its rotational source is Haar and
carries no orientational information whatsoever, and indeed the Haar measure
has no mean to centre. Section~\ref{sec:future-source} reports a measurement of
that asymmetry and a negative result on one attempt to remove it.

% ===========================================================================
\chapter{SigmaFlow-Minimal: The Minimal Conversion}
\label{ch:minimal}

\section{Design principle and what actually changed}
\label{sec:minimal-change}

The comparison is only interpretable if the substitution is confined to the
generative mechanism. This was treated as a constraint on the implementation
rather than as an aspiration, and the resulting partition is given in
Table~\ref{tab:minimal-change}.

\begin{table}[htbp]
\centering
\caption{Partition of the system under the minimal-change constraint.
Components marked identical were verified by file hash over the relevant
source trees.}
\label{tab:minimal-change}
\small
\begin{tabular}{@{}p{0.30\textwidth} p{0.16\textwidth} p{0.46\textwidth}@{}}
\toprule
\textbf{Component} & \textbf{Status} & \textbf{Note} \\
\midrule
EquiformerV2 backbone            & identical & all backbone source files hash-identical between arms \\
Featurisation, graph construction & identical & inherited \\
Data pipeline, fragmentation      & identical & same conformers, same splits \\
Output block ($\ell=1$ head)      & identical & one field $f_i$, Eq.~\eqref{eq:force-field} \\
Newton-Euler reduction            & identical & Eqs.~\eqref{eq:newton-euler}--\eqref{eq:updates} \\
Optimiser and schedule            & identical & AdamW \citep{loshchilov2019adamw}, same cosine schedule \\
Loss weights                      & identical & SigmaDock production values \\
\midrule
Forward noising process           & \textbf{removed} & replaced by an explicit probability path \\
Score rescaling by $\lambda(t)$   & \textbf{removed} & the scalar apparatus converting scores to displacements is deleted, not modified \\
Probability path                  & \textbf{new} & Eqs.~\eqref{eq:sf-trans}--\eqref{eq:sf-rot} \\
Training target                   & \textbf{new} & conditional velocity in place of score \\
Frame conversion                  & \textbf{new} & Eq.~\eqref{eq:frame-conversion} \\
Sampler                           & \textbf{new} & deterministic Euler, Eq.~\eqref{eq:euler} \\
\bottomrule
\end{tabular}
\end{table}

Two residual confounders should be recorded rather than glossed. The two arms
run in separate software environments with different interpreter versions, and
the SigmaDock arm is invoked with its own rotational-score configuration flags,
which have no counterpart in a flow-matching model. Neither touches the
backbone, but neither is strictly nothing.

\section{The vector field readout}
\label{sec:sf-readout}

The network output and its Newton-Euler reduction are inherited exactly, so the
per-fragment increments $\Delta T_F$ and $\skewm{\omega}_F$ of
\eqref{eq:updates} are available unchanged. What differs is their
interpretation. SigmaFlow reads them as a velocity, which requires no
time-dependent rescaling, so the translational prediction is
\begin{equation}
  \pred^{\mathrm{trans}} \;=\; \Delta T_F
  \;=\; \frac{1}{|F|}\sum_{i \in F} f_i .
  \label{eq:sf-trans-pred}
\end{equation}

The rotational prediction requires one further step, which is the only piece of
new mathematics in the minimal conversion. As noted in
Section~\ref{sec:trivialisation}, $\skewm{\omega}_F$ is a spatial-frame
quantity, whereas the target \eqref{eq:sf-rot} and the integrator
\eqref{eq:euler} are both body-frame. Requiring the two descriptions of the
same infinitesimal motion to agree,
\begin{equation}
  \Rot_t \Exp\!\bigl(\skewm{\Omega}\,\Delta t\bigr)
  \;=\; \Exp\!\bigl(\skewm{w}\,\Delta t\bigr) \Rot_t ,
\end{equation}
gives the adjoint transport of \eqref{eq:adjoint}, so the predicted rotational
velocity is
\begin{equation}
  \pred^{\mathrm{rot}}
  \;=\; \Rot_t^{\top}\, \skewm{\omega}_F\, \Rot_t .
  \label{eq:frame-conversion}
\end{equation}
Conjugating a skew-symmetric matrix by a rotation preserves skew-symmetry, so
$\pred^{\mathrm{rot}}$ remains a valid element of $\sothree$ and every
downstream shape is unchanged.

This conversion is required by the construction rather than by any particular
readout. The target $\Log(\Rot_t^\top \Rot_1)$ is \emph{invariant} under a
global rotation $Q$, since $(Q\Rot_t)^\top(Q\Rot_1) = \Rot_t^\top\Rot_1$,
whereas any equivariant network output transforms covariantly. The conjugation
is what reconciles the two, and it would be needed for any equivariant
rotational readout, not only for the Newton-Euler one. The translational
channel needs no counterpart, because $\Delta T_F$ and
$\vfield^{\mathrm{trans}}_t$ are both expressed in the world frame.

\section{Deterministic integration}
\label{sec:integration}

Sampling integrates \eqref{eq:euler} from $t=0$ to $t=1$ over a fixed number of
steps. On the translational factor this is ordinary explicit Euler. On the
rotational factor the update is a retraction, meaning the step is taken along
the manifold rather than in an ambient space followed by a projection, and
because $\Exp$ maps into $\SOthree$ exactly, every iterate is a rotation matrix
by construction and no re-orthogonalisation is required.

Three properties distinguish this from the reverse-time stochastic differential
equation it replaces. Sampling is deterministic, so the generated pose is a
function of the source draw alone and the diversity of a $K$-sample set is
inherited entirely from $p_0$. The sampler has no temperature or noise scale to
tune. And the number of network evaluations equals the number of steps and is a
free parameter at inference, decoupled from anything fixed at training time.
Because the conditional paths are straight in $\R^3$ and geodesic on
$\SOthree$, the marginal field they induce is expected to be nearly straight,
so few Euler steps should suffice. That is a prediction of the construction
rather than a hope, and Section~\ref{sec:minimal-nfe} tests it. The step count
inherited from SigmaDock is $25$, a value tuned there for a stochastic reverse
process, and there is no reason for that number to transfer.

\section{Results}
\label{sec:minimal-results}

\begin{reserved}{SigmaFlow-Minimal against SigmaDock at matched budget}
\textbf{Experiment.} The $72$\,h matched pair, SigmaFlow-Minimal against
SigmaDock, identical data, identical wall-clock allocation on identical
hardware, identical batch size, optimiser and annealing schedule. Sampling with
$K=10$ independent seeds per complex over the $209$ evaluation complexes.\\
\textbf{Metrics.} Top-1 and oracle-at-$K$ rates under \eqref{eq:success}, the
RMSD distribution, PoseBusters validity, RMSD after rigid superposition and
centroid displacement.\\
\textbf{Question answered.} Whether replacing score-based diffusion by
Riemannian flow matching changes docking accuracy at matched compute, and in
which direction.\\
\textbf{Dependence.} If the two arms are within sampling error, the thesis
reports a null result on the mechanism and the interest shifts to the
inference-time properties in Sections~\ref{sec:minimal-nfe} and
\ref{sec:minimal-seeds}. If they differ, the diagnostic quantities in
Section~\ref{sec:minimal-rotation} are the first place to look for the reason.
\end{reserved}

\figplaceholder{fig:minimal-curves}
{Evolution of the pose-level metrics over training for SigmaFlow-Minimal and
SigmaDock at matched budget.}
{Three stacked panels sharing an $x$ axis in optimisation steps, with a second
axis in wall-clock hours. Panel 1: median RMSD. Panel 2: fraction with
RMSD $<2$\,\AA. Panel 3: fraction PoseBusters-valid, and the conjunction
\eqref{eq:success} as a separate line. Two colours, one per arm, with bootstrap
bands over complexes. Snapshots from a single annealing schedule are drawn as
curves; separately annealed endpoints are drawn as isolated markers and are
never joined to the curves, because a mid-anneal snapshot is not the model one
would obtain by training to that budget with its own schedule.}

\figplaceholder{fig:minimal-pymol}
{Generated fragment trajectories for SigmaFlow-Minimal on representative
complexes, rendered in PyMOL.}
{A grid of panels, one per selected complex, chosen to span the fragment-count
range of Figure~\ref{fig:fragment-histogram}, for instance one ligand each with
$N = 1$, $3$, $5$ and $8$. In each panel show the protein pocket surface in
light grey, the crystallographic pose in a saturated reference colour, and the
generated pose. Overlay the integration trajectory by drawing the fragment
centroids at every integration step as a fading trace, with each fragment in
its own colour so that the reader can see whether fragments move coherently or
independently. State the number of integration steps in the caption. The point
of the figure is qualitative, namely whether failures look like a
mis-positioned but coherent ligand or like a set of fragments that never
assembled.}

\subsection{Number of function evaluations}
\label{sec:minimal-nfe}

\begin{reserved}{Accuracy against integration steps}
\textbf{Experiment.} Sample the same trained checkpoint with
$\{1, 5, 10, 25, 50, 100\}$ Euler steps, fixed seeds, everything else held
constant.\\
\textbf{Metrics.} The criteria of \eqref{eq:success} against step count,
together with wall-clock time per pose.\\
\textbf{Question answered.} Whether the near-straight conditional paths
predicted in Section~\ref{sec:integration} translate into accurate sampling at
low step counts, and where the accuracy-cost curve saturates.\\
\textbf{Dependence.} A curve that saturates well below $25$ steps is a concrete
inference-time advantage of the flow formulation and belongs in the abstract. A
curve that does not saturate until far above $25$ would falsify the
straightness argument and require explanation.\\
\textbf{Caveat already established.} Wall-clock timings from the preliminary
sweep are not usable, because they show a spurious minimum at intermediate step
counts that is the signature of a cold filesystem cache on the first array
task. The step-count axis itself is unaffected.
\end{reserved}

\figplaceholder{fig:minimal-nfe}
{Docking accuracy as a function of the number of integration steps.}
{Accuracy on the $y$ axis, number of Euler steps on a logarithmic $x$ axis,
with one line per metric in \eqref{eq:success} and a bootstrap band. Mark the
inherited default of $25$ steps with a vertical rule. If a second axis is used
for wall-clock cost per pose, take it from a clean timing run rather than from
the preliminary sweep.}

\subsection{Number of samples}
\label{sec:minimal-seeds}

\begin{reserved}{Accuracy against the number of independent samples}
\textbf{Experiment.} For $K \in \{1,\ldots,10\}$, report both the oracle rate
over $K$ draws and the Top-1 rate after the inherited ranking heuristic.\\
\textbf{Metrics.} Oracle-at-$K$ and Top-1 under \eqref{eq:success}, and the
ratio between them.\\
\textbf{Question answered.} How much of the attainable accuracy is lost in
ranking rather than in generation, which is the quantity
Chapter~\ref{ch:confidence} sets out to recover.\\
\textbf{Pre-registered.} The ratio of oracle-at-$10$ to oracle-at-$1$ was fixed
as a reported quantity before the long runs, together with the interpretation
of each outcome, and is not to be reinterpreted afterwards. On the preliminary
$12$\,h pair the ratio was $6.8$ for SigmaFlow and $4.4$ for SigmaDock.
\end{reserved}

\figplaceholder{fig:minimal-seeds}
{Oracle and Top-1 accuracy as a function of the number of independent samples.}
{Two lines per arm, oracle-at-$K$ and Top-1 after ranking, with $K$ on the $x$
axis. Shade the region between them, since that area is precisely the accuracy
that a better ranker could recover and is the motivation for
Chapter~\ref{ch:confidence}.}

\subsection{The rotational channel}
\label{sec:minimal-rotation}

The diagnostic quantity of most interest is the per-fragment rotational error
against the null of \eqref{eq:haar-null}, because it separates a model that
places fragments correctly but orients them poorly from one that fails at
both.

\begin{reserved}{Rotational error against the Haar null}
\textbf{Experiment.} Per-fragment geodesic angle between the generated and the
crystallographic fragment orientation, aggregated over the evaluation set, for
both arms and at several points along training.\\
\textbf{Metrics.} Median and mean geodesic angle against the null values
$\Ang{132.3}$ and $\Ang{126.5}$, reported with bootstrap intervals over
complexes.\\
\textbf{Question answered.} Whether the rotational channel carries signal at
this budget, and whether any deficit narrows as training proceeds.\\
\textbf{Deliberately not developed further here.} On the preliminary $6$\,h run
the absolute rotational error did not fall below the null, but it is not yet
known whether that persists under longer training. This subsection therefore
reports the measurement and stops. Any mechanistic account would be premature,
and if the deficit does persist at full budget it belongs in
Section~\ref{sec:discussion-open} as an open question rather than as a claim.
\end{reserved}

% ===========================================================================
\chapter{SigmaFlow-Separate: Translation and Rotation Specific Heads}
\label{ch:separate}

\section{Motivation}

The minimal conversion inherits a readout in which one per-atom field serves
both channels. As established in Section~\ref{sec:backbone}, translation reads
the zeroth spatial moment of $f$ over a fragment and rotation reads its first
moment about the fragment centre, and the centring identity
$\sum_{i\in F}(r_i - c_F) = 0$ means that these two moments are not merely
different summaries but pick out orthogonal spatial modes. A spatially constant
component of the field is invisible to the rotational channel, and a
zero-sum spatially varying pattern is invisible to the translational one.

Under diffusion this coupling is mitigated by the fact that the two channels
are rescaled independently by their respective noise schedules, so the
effective gain on each is set separately at every noise level. Under flow
matching that apparatus is deleted, as recorded in
Table~\ref{tab:minimal-change}, and the two channels read a single unscaled
field. Whether the single field can serve both is then an empirical question,
and it is the question this chapter isolates.

A second, narrower motivation concerns conditioning. The rotational path passes
through $I_F^{-1}$, and the inertia tensor of a fragment consisting of two
atoms is that of a line, which is singular in the direction of the bond. The
implementation regularises the eigenvalues before inversion, but the
regularised tensor of a two-atom fragment still has a condition number of order
$10^{8}$, so the rotational channel is numerically delicate for exactly those
fragments that the fragmentation scheme produces most readily. A readout that
does not invert an inertia tensor avoids the issue entirely.

\begin{remark}
An earlier version of this argument appealed to single-atom fragments, for
which the torque vanishes identically. That argument does not apply, because
the baseline asserts a minimum fragment size of two atoms and single-atom
fragments therefore never arise. The two-atom case above is the one that
survives scrutiny.
\end{remark}

\section{Construction}

The variant replaces the single output block by two blocks of identical
configuration on a shared trunk, so that the backbone up to the final
attention block is untouched and only the readout is duplicated. Each block
emits its own $\ell=1$ field per atom,
\begin{equation}
  f_i^{T} \in \R^3, \qquad f_i^{R} \in \R^3 .
\end{equation}
The Newton-Euler reduction is then removed from the active path and replaced by
mean pooling over the atoms of each fragment,
\begin{equation}
  \pred^{\mathrm{trans}} \;=\; \frac{1}{|F|}\sum_{i\in F} f_i^{T},
  \qquad
  \bar{w}_F \;=\; \frac{1}{|F|}\sum_{i\in F} f_i^{R},
  \qquad
  \pred^{\mathrm{rot}} \;=\; \Rot_t^{\top}\,\skewm{\bar{w}_F}\,\Rot_t .
  \label{eq:separate-readout}
\end{equation}
The translational expression is unchanged in form, since uniform masses already
made \eqref{eq:sf-trans-pred} a mean. The rotational expression differs
substantively, since the torque and the inverse inertia are gone and the
network is asked to emit an angular velocity directly. The frame conversion of
\eqref{eq:frame-conversion} is retained without modification, for the reason
given in Section~\ref{sec:sf-readout}, namely that it is required by the
invariance of the target rather than by the particular readout.

Mean pooling rather than summation is deliberate. A sum makes the predicted
velocity scale with fragment size, which over a benchmark whose fragments range
from two atoms to several dozen would impose an uncontrolled size-dependent
gain on the loss. A mean is size-invariant.

The cost is one additional attention block. Measured on the production
configuration the model has $16{,}899{,}826$ parameters against
$14{,}979{,}377$ for the single-head model, of which $1{,}920{,}449$ belong to
each of the two identically configured heads, an increase of $12.82\%$.
Everything else, including the data pipeline, the fragmentation, the optimiser,
the schedule, the loss weights, the source distribution and the sampler, is
byte-identical to SigmaFlow-Minimal.

The construction has one structural property worth recording, because it is not
available in the single-head model and it is what the variant is for. With two
heads the translational loss reaches the rotational block with a gradient that
is exactly zero and the rotational loss reaches the translational block with a
gradient that is exactly zero, while both losses reach the shared trunk. The
two channels are therefore structurally decoupled at the readout and coupled
only through the representation. In the single-head model no such statement is
possible, since both losses share one output.

\section{Results}

\begin{reserved}{SigmaFlow-Separate against SigmaFlow-Minimal}
\textbf{Experiment.} A run at wall-clock budget matched to the corresponding
SigmaFlow-Minimal run, identical in every other respect. Compute rather than
step count is the matched quantity, so the larger model completes fewer
optimisation steps within the same allocation, and the achieved step count is
reported alongside every comparison.\\
\textbf{Metrics.} The pose-level criteria of \eqref{eq:success}, the
translational and rotational validation losses reported separately, and the
per-fragment rotational error against the null of \eqref{eq:haar-null}.\\
\textbf{Question answered.} Whether the shared readout was a binding constraint
on the rotational channel, or whether the deficit lies elsewhere.\\
\textbf{Dependence.} An improvement concentrated in the rotational channel
supports the moment-competition account of Section~\ref{sec:backbone}. An
improvement in neither channel, or one attributable purely to the extra
capacity, argues that the readout was not the limiting factor and redirects
attention to the source distribution of Section~\ref{sec:future-source}.\\
\textbf{Control for capacity.} The variant adds parameters as well as
separating the channels, so an unambiguous attribution would require a
capacity-matched control with a widened single head. That control is recorded
here as a limitation rather than run.
\end{reserved}

\figplaceholder{fig:separate-curves}
{Metric evolution over training for SigmaFlow-Separate against
SigmaFlow-Minimal at matched compute.}
{Same three-panel layout as Figure~\ref{fig:minimal-curves}, with the two
SigmaFlow variants as the two series and SigmaDock retained as a faint
reference line. Because the budget is matched in compute rather than in steps,
plot against wall-clock hours as the primary axis and mark the achieved step
count of each arm at its endpoint.}

\figplaceholder{fig:separate-pymol}
{Generated fragment trajectories for SigmaFlow-Separate on the same complexes
as Figure~\ref{fig:minimal-pymol}.}
{Identical rendering protocol, camera and colour scheme to
Figure~\ref{fig:minimal-pymol}, on the same complexes, so that the two figures
can be compared panel by panel. Any difference in orientation behaviour should
be visible without consulting the numbers.}

\figplaceholder{fig:separate-nfe-seeds}
{Integration steps and sample count for SigmaFlow-Separate.}
{Two panels reproducing the protocols of Figures~\ref{fig:minimal-nfe} and
\ref{fig:minimal-seeds} for this variant, so that the inference-time
characterisation is available for every model in the thesis rather than for the
baseline alone.}

% ===========================================================================
\chapter{SigmaFlow-Confidence: An Integrated Confidence Model}
\label{ch:confidence}

\section{Why ranking is a separate problem}

Every generative docking method samples repeatedly and must then choose. The
quantity reported as accuracy therefore depends on a ranker as much as on a
generator, and the two can fail independently. Section~\ref{sec:minimal-seeds}
measures the size of that gap directly, and on the preliminary pair it was
large, with the oracle rate over ten samples exceeding the single-sample rate
by a factor of $6.8$ for SigmaFlow and $4.4$ for SigmaDock. A ranker that
recovered even part of that difference would be worth more than a
correspondingly sized improvement in the generator.

The two systems this thesis builds on resolve the problem differently.
DiffDock \citep{corso2023diffdock} trains a separate \emph{confidence model}.
The generative model is first trained and then frozen, and is used to sample
poses for the training complexes. Each sampled pose is labelled by whether its
RMSD to the crystallographic pose falls below the success threshold, and a
second, smaller equivariant network is trained on those labelled samples as a
binary classifier. At inference the generator produces $K$ poses and the
confidence model ranks them. The essential point, and the reason the confidence
model is not simply trained on crystal poses, is that the distribution it must
discriminate over is the generator's own output distribution, which is not the
data distribution. A classifier trained on crystallographic poses against
arbitrary decoys would be solving a different problem from the one it faces at
inference.

SigmaDock deliberately declines this. It reports that the reliability of its
samples makes a learned filter unnecessary and ranks instead by a fixed
heuristic combining a Vinardo pseudo binding energy
\citep{quiroga2016vinardo} with a set of physicochemical checks
\citep{prat2026sigmadock}. Its ablations show that the heuristic is doing real
work rather than being a formality, since removing the energy term or the
physicochemical checks each cost several points of Top-1 accuracy.

This leaves an open question that the present setting is well placed to
examine. SigmaDock's argument for the heuristic rests on the quality of its
samples at full training budget. At the budgets available here the samples are
considerably worse, and the oracle-to-Top-1 gap is correspondingly larger, so
the regime in which a learned confidence model helps is precisely the regime
this thesis operates in.

\section{Construction}

The variant follows DiffDock's protocol, adapted to the fragment
representation and reusing the inherited backbone so that the addition remains
a controlled one.

\begin{enumerate}[leftmargin=*,itemsep=3pt]
\item Train SigmaFlow to convergence under the budget and freeze it.
\item For each training complex, sample $K$ poses from the frozen generator
      using the deterministic sampler of \eqref{eq:euler} with independent
      source draws, so that the sampled distribution is exactly the one seen at
      inference.
\item Label each pose by the criterion of \eqref{eq:success}, giving a binary
      target per pose.
\item Train a confidence head $d_\psi(x,\ctx) \in \R$ on these labelled samples
      with the binary cross-entropy objective
      \begin{equation}
        \Loss_{\mathrm{conf}}(\psi)
        = -\E\bigl[\, y \log \varsigma(d_\psi) + (1-y)\log(1-\varsigma(d_\psi))\,\bigr],
        \label{eq:conf-loss}
      \end{equation}
      where $\varsigma$ is the logistic function and $y$ the label. The head
      reads the same equivariant features as the generator but pools to a
      single invariant scalar, since a confidence score must be invariant
      rather than equivariant.
\item At inference, sample $K$ poses and return the one maximising $d_\psi$.
\end{enumerate}

Two design points require care and are recorded as such. The labels are highly
imbalanced at low training budgets, because most sampled poses fail, so the
positive rate must be reported and the objective weighted or the sampling
stratified accordingly. And the confidence model must be trained on samples
from complexes in the training split only, since sampling from evaluation
complexes to train the ranker would leak the answer into the reported metric.

\section{Results}

\begin{reserved}{Learned ranking against the inherited heuristic}
\textbf{Experiment.} With the generator fixed, compare three rankers on the
same set of $K$ generated poses per complex, namely the inherited
energy-plus-checks heuristic, the learned confidence model, and the oracle.\\
\textbf{Metrics.} Top-1 under \eqref{eq:success} for each ranker, the fraction
of the oracle-to-heuristic gap that the learned model recovers, and the
calibration of $\varsigma(d_\psi)$ against the empirical success rate.\\
\textbf{Question answered.} Whether a learned confidence model is worth its
training cost at this budget, and whether SigmaDock's decision to omit one is
budget-dependent.\\
\textbf{Dependence.} Because the generator is held fixed, this comparison is
clean regardless of how the generative comparison in
Chapter~\ref{ch:minimal} turns out, and it remains reportable even if the
matched-budget result there is null.
\end{reserved}

\figplaceholder{fig:confidence-ranking}
{Top-1 accuracy under three rankers, against the number of generated samples.}
{Three lines against $K$, namely the inherited heuristic, the learned
confidence model, and the oracle, all on the same generated pose sets so that
the generator is held fixed. Shade the region recovered by the learned model as
a fraction of the region between heuristic and oracle, since that fraction is
the headline quantity of the chapter.}

\figplaceholder{fig:confidence-calibration}
{Calibration of the learned confidence score.}
{A reliability diagram with predicted confidence on the $x$ axis, binned, and
the empirical fraction of poses satisfying \eqref{eq:success} on the $y$ axis,
with the diagonal drawn. Include a histogram of predicted confidences beneath,
since calibration in a sparsely populated bin is not informative. Report the
positive rate of the training labels in the caption, because the imbalance
noted above determines how this figure should be read.}

\figplaceholder{fig:confidence-pymol}
{Poses selected by the heuristic and by the learned confidence model on
complexes where the two disagree.}
{Paired panels for a handful of complexes on which the two rankers select
different poses. Show the crystallographic pose, the heuristic's choice and the
confidence model's choice in three distinguishable colours, with the RMSD of
each annotated. Choose examples in both directions, that is, cases where the
learned model wins and cases where it loses, rather than only favourable ones.}

% ===========================================================================
\chapter{Discussion}
\label{ch:discussion}

\section{What the comparison establishes}

\begin{reserved}{Synthesis across the three variants}
To be written once Chapters~\ref{ch:minimal}, \ref{ch:separate} and
\ref{ch:confidence} carry their results. The section should state, in order,
what the matched-budget comparison shows about the generative mechanism, what
the separated readout adds, and how much of the remaining gap is attributable
to ranking rather than to generation. It should also state plainly which of the
three questions the available compute was not sufficient to answer.
\end{reserved}

\section{Limitations}

Several limitations constrain what can be claimed, and they are stated here
rather than distributed through the results.

The compute budget is the binding one. SigmaDock's published results were
obtained with $256$ epochs of training, whereas the runs reported here reach a
small fraction of that. Absolute accuracies are correspondingly far below the
literature for every arm, including the reproduced baseline, and no benchmark
claim is made. What the design supports is a comparison between mechanisms at
equal cost, and the reader should resist reading any absolute number here
against a published one.

The evaluation set differs from the published one, and the difference is not
fully characterised. Section~\ref{sec:evaluation-set} records what is known,
namely that the copy of the benchmark available to this project contains $209$
complexes rather than the $308$ the paper evaluates, and that the reduction was
not reconstructible. Whether the $209$ are a subset of the published $308$ has
not been established. Every absolute number in this thesis should therefore be
read as pertaining to this particular $209$-complex set, and the internal
matched comparison, which uses the same set for every arm, is unaffected.

The comparison in Chapter~\ref{ch:separate} matches compute rather than
optimisation steps, which is the correct control for a question about
mechanisms at fixed cost but means the larger model takes fewer steps. The
achieved step count is therefore reported alongside every such comparison.
Further, that variant adds capacity at the same time as it separates the
channels, and the capacity-matched control was not run.

Snapshots taken during a single annealing schedule are not equivalent to
separately annealed shorter runs, because a mid-schedule snapshot sits at a
learning rate far above its terminal value. Training trajectories are therefore
plotted as curves under a fixed shared schedule, separately annealed endpoints
are plotted as isolated markers, and the two are never joined. No claim of the
form ``performance scales with compute as $\ldots$'' is made, since
establishing that would require a fully annealed run per budget.

Finally, one residual environmental confounder remains between the two arms,
since they run under different interpreters and the baseline is invoked with
configuration flags that have no flow-matching counterpart. Neither touches the
backbone, and both are recorded rather than corrected.

\section{Open questions}
\label{sec:discussion-open}

\begin{reserved}{Open questions arising from the results}
To be written with the results. The rotational channel of
Section~\ref{sec:minimal-rotation} is the natural candidate, provided the
measurement at full budget supports raising it. If the deficit against the Haar
null has closed by then, this section should say so and the question is
retired.
\end{reserved}

\section{Future work}

\subsection{Informative source distributions}
\label{sec:future-source}

Section~\ref{sec:source-theory} established that the source distribution is a
free modelling choice under flow matching and a fixed consequence of the
forward process under diffusion. This is the most distinctively flow-specific
degree of freedom available, and it is unexploited in the present work, since
SigmaFlow inherits \eqref{eq:stationary} unchanged.

The asymmetry between the two factors is what makes the idea attractive and can
be quantified. On the translational factor the source is already informative,
because pocket-centred coordinates place its mean near the answer. Measured
over the evaluation set, the median normalised distance from a standard normal
draw to the true fragment centre is $2.1$, against $1.5$ for a draw centred on
the pocket centre. On the rotational factor there is no counterpart, because
the Haar measure has no mean and the median geodesic distance from a Haar draw
to the true orientation is $\Ang{131.3}$, essentially the null of
\eqref{eq:haar-null}.

A natural conjecture is that a context-dependent rotational source
$p_0(\cdot\mid\ctx)$ concentrated near plausible orientations would shorten the
transport paths and simplify the learning problem. One family of such
heuristics was tested and \emph{failed}, and the result is reported here
because a negative result obtained under a pre-registered gate is worth more
than an untested conjecture. Aligning the principal axes of the fragment
conformer with those of the pocket gave a median geodesic distance of
$\Ang{138.1}$, which is $\Ang{6.7}$ \emph{worse} than the uninformed Haar
source. A control that placed $\Rot_0$ at the identity gave $\Ang{130.7}$, and
a leakage control in which the heuristic was allowed to see the
crystallographic geometry still gave only $\Ang{135.4}$, so the ceiling of this
family lies below the uninformed source. The measurement itself is sound, since
the Haar arm reproduced the analytic null of \eqref{eq:haar-null} to within
about one degree.

The scope of that conclusion should be stated narrowly. One heuristic family
was tested, so what is established is that principal-axis conditioning does not
help, not that no informative rotational source can help. Directions that
remain open include mixtures on $\SOthree$ that hedge the axis ambiguity
inherent in aligning near-symmetric fragments, sources that couple the
fragments rather than treating them independently, and a source learned jointly
with the flow, the last of which would require care to remain leakage-free
since a learned source has the capacity to memorise the answer.

\subsection{Likelihood-based pose confidence}
\label{sec:future-likelihood}

Chapter~\ref{ch:confidence} obtains a confidence score by training a
classifier. The flow formulation offers a second route that the diffusion
formulation does not offer as directly, namely reading a likelihood off the
generative process itself.

Because sampling integrates a deterministic ordinary differential equation, the
model is a continuous normalising flow \citep{chen2018neural}, and the density
of a generated point follows from the instantaneous change of variables. On
$\R^{d}$,
\begin{equation}
  \log p_1(x_1) \;=\; \log p_0(x_0)
    \;-\; \int_0^1 \divg \pred(t, x_t)\, \dif t ,
  \label{eq:cnf-likelihood}
\end{equation}
where $x_t$ solves \eqref{eq:flow-ode}. The same statement holds on a
Riemannian manifold with the divergence taken with respect to the Riemannian
volume form \citep{chen2024riemannian}, which on $\SOthree$ with the
bi-invariant metric is the Haar measure, so the construction is available on
$\Mfd$ without further choices. Evaluating \eqref{eq:cnf-likelihood} for a
\emph{given} pose requires integrating the flow backwards from that pose to
$t=0$, which is a well-posed problem precisely because the dynamics are
deterministic and reversible.

The practical obstacle is the divergence term, whose exact evaluation costs
$O(d)$ backward passes. The standard remedy is the Hutchinson trace estimator
\citep{hutchinson1989stochastic,grathwohl2019ffjord},
\begin{equation}
  \divg \pred \;=\; \E_{\epsilon}\Bigl[\epsilon^{\top}
    \frac{\partial \pred}{\partial x}\, \epsilon\Bigr],
  \qquad \E[\epsilon] = 0, \quad \operatorname{Cov}(\epsilon) = I,
  \label{eq:hutchinson}
\end{equation}
which replaces the exact trace by a stochastic estimate at the cost of one
vector-Jacobian product per sample. On $\Mfd$ the estimator must draw
$\epsilon$ from the tangent space at $x_t$ with the metric-induced covariance,
which on the rotational factors means sampling in $\sothree$ under the
bi-invariant inner product.

Three uses suggest themselves, in increasing order of ambition. The conditional
log-density $\log p_\theta(x\mid\ctx)$ could rank the $K$ generated poses
directly, giving a confidence score that requires no additional training and no
labelled samples, and which could be compared against the learned classifier of
Chapter~\ref{ch:confidence} on identical pose sets. The same quantity evaluated
at the crystallographic pose would give a likelihood-based evaluation metric
that is independent of RMSD thresholds. And the spread of $\log p_\theta$ over a
sample set would give a measure of the model's own uncertainty for a given
complex, which is the quantity classical docking cannot provide and which
motivated the generative formulation in the first place.

Two cautions belong with this proposal. A high model density is a statement
about the model, not about binding affinity, so calibration against measured
outcomes would be required before any such score could be read as a physical
quantity. And the estimator in \eqref{eq:hutchinson} is unbiased for the
divergence but not for the log-density after integration, so the variance of
the resulting estimate must be reported rather than assumed small.

\subsection{Further directions}

Three smaller directions are worth recording. The integration scheme is
explicit Euler throughout, chosen for correctness rather than efficiency, and
the near-straight paths of Section~\ref{sec:integration} suggest that a
higher-order or adaptive scheme would either reduce the step count further or
show that the paths are straight enough that it cannot. The product structure
of $\Mfd$ is imposed rather than derived, and a construction that coupled the
rotational and translational fields, for instance through the semidirect
product structure of $\SEthree$, would be a genuine change of model rather than
of parameterisation. Finally, SigmaDock notes that its fragment formulation
extends to flexible docking by treating selected side chains as additional
fragments, and nothing in the flow-matching substitution obstructs that
extension.

% ===========================================================================
\chapter{Conclusion}
\label{ch:conclusion}

\begin{reserved}{Conclusion}
To be written last. It should restate what was substituted and what was held
fixed, give the verdict of the matched-budget comparison in one sentence,
state what the two variants added, and close on the two directions of
Sections~\ref{sec:future-source} and \ref{sec:future-likelihood} that the flow
formulation makes available and the diffusion formulation does not. It should
not introduce any quantity that has not appeared earlier.
\end{reserved}

% ---------------------------------------------------------------------------
\bibliographystyle{plainnat}
\bibliography{../refs}

% ===========================================================================
\appendix

\chapter{Equivariance and the Inherited Backbone}
\label{app:equivariance}
\dnote{DRAFT}{Retained from the earlier draft and not currently referenced by
the main text. The main text is written to stand without the appendices.}

\chapter{Diffusion Background}
\label{app:diffusion}
\dnote{DRAFT}{Retained from the earlier draft.}

\chapter{Derivations}
\label{app:derivations}
\dnote{DRAFT}{Retained from the earlier draft.}

\chapter{Numerical Verification}
\label{app:numerics}
\dnote{DRAFT}{Retained from the earlier draft.}

\chapter{Extended Results}
\label{app:extended}
\dnote{DRAFT}{Retained from the earlier draft.}

\chapter{Implementation Reference}
\label{app:implementation}
\dnote{DRAFT}{Retained from the earlier draft.}

\end{document}
